\documentclass[12pt, a4paper]{article}

% -------------------------------------------------------
% Packages
% -------------------------------------------------------
\usepackage[a4paper, margin=1in]{geometry}
\usepackage{graphicx}
\usepackage{booktabs}
\usepackage{array}
\usepackage{xcolor}
\usepackage{hyperref}
\usepackage[T1]{fontenc}
\usepackage[utf8]{inputenc}
\usepackage{fancyhdr}
\usepackage{tabularx}
\usepackage{longtable}
\usepackage{enumitem}
\usepackage{titlesec}
\usepackage{parskip}
\usepackage{svg}

% -------------------------------------------------------
% Hyperlink styling
% -------------------------------------------------------
\hypersetup{
    colorlinks=true,
    linkcolor=black,
    urlcolor=blue,
    citecolor=black,
    pdftitle={Beacon SRS},
    pdfauthor={Rudra Ponkshe},
    pdfsubject={Software Requirements Specification}
}

% -------------------------------------------------------
% Header / Footer
% -------------------------------------------------------
\pagestyle{fancy}
\fancyhf{}
\rhead{\small SRS Document}
\lhead{\small Beacon --- Location Based Attendance Tracking System}
\rfoot{\small Page \thepage}
\lfoot{\small \href{https://gitlab.com/realRudraP/beacon}{gitlab.com/realRudraP/beacon}}
\renewcommand{\headrulewidth}{0.4pt}
\renewcommand{\footrulewidth}{0.4pt}

% -------------------------------------------------------
% Section title formatting
% -------------------------------------------------------
\titleformat{\section}{\large\bfseries}{{\thesection}}{1em}{}[\titlerule]
\titleformat{\subsection}{\normalsize\bfseries}{\thesubsection}{1em}{}
\titleformat{\subsubsection}{\normalsize\itshape\bfseries}{\thesubsubsection}{1em}{}

% -------------------------------------------------------
% Useful shorthand for FR/NFR blocks
% -------------------------------------------------------
\newcommand{\frblock}[2]{%
    \noindent\textbf{#1}\hfill\textit{Priority: #2}%
    \par\smallskip%
}
\newcommand{\srsfield}[1]{\noindent\textbf{#1:}\par}

% ================================================================
% DOCUMENT
% ================================================================
\begin{document}

% -------------------------------------------------------
% COVER PAGE
% -------------------------------------------------------
\begin{titlepage}
    \centering
    \thispagestyle{empty}

    \noindent\rule{\linewidth}{1.5pt}
    \vspace{0.5em}

    {\large\scshape Software Requirements Specification}

    \vspace{0.3em}
    \noindent\rule{\linewidth}{0.5pt}
    \vspace{2em}

    {\Huge\bfseries Beacon}

    \vspace{0.5em}
    {\large Location Based Attendance Tracking System}\\[4pt]
    {\normalsize Offline-First, Self-Hostable, Open Source}

    \vspace{2em}

    \fbox{%
        \parbox{0.25\linewidth}{%
            \centering
            \textbf{Version}\\[4pt]
            1.2 --- Draft
        }%
    }

    \vfill

    \begin{tabularx}{0.85\linewidth}{@{} l X @{}}
        \toprule
        \textbf{Author}       & Rudra Ponkshe \\
        \midrule
        \textbf{Repository}   & \url{https://gitlab.com/realRudraP/beacon} \\
        \midrule
        \textbf{Licence}      & To be determined \\
        \midrule
        \textbf{Date}         & April 2026 \\
        \bottomrule
    \end{tabularx}

    \vspace{2em}
    \noindent\rule{\linewidth}{1.5pt}
\end{titlepage}

% -------------------------------------------------------
% REVISION HISTORY
% -------------------------------------------------------
\newpage
\thispagestyle{fancy}
\section*{Revision History}
\addcontentsline{toc}{section}{Revision History}

\begin{tabularx}{\linewidth}{@{} c c l X @{}}
    \toprule
    \textbf{Version} & \textbf{Date} & \textbf{Author} & \textbf{Description} \\
    \midrule
    0.1 & 08 April 2026 & Rudra Ponkshe & Initial draft setup \\
    1.0 & 09 April 2026 & Rudra Ponkshe & First complete SRS version for TA-2 submission \\
    1.1 & 12 April 2026 & Rudra Ponkshe & Added: worker attendance history, wage calculation,
                                     remote system configuration, India public holidays,
                                     per-worker weekly off, absentia marking, petty-cash
                                     ledger, clock-in reminder notifications, short-hours
                                     warning, and under-hours dashboard flag. \\
    1.2 & 17 April 2026 & Rudra Ponkshe & Open source revision. The project is now set to
    natively support multi-tenant deployments, the India specific defaults are now configurable as locale-configurable records. Pluggable wage collection strategies are introduced. Push Notifications and crash reporting providers are introduced as optional and swappable. Added open source audience and reframed out of scope items as community contribution roadmap\\
    1.3 & 18 April 2026 & Rudra Ponkshe & \textbf{ADR-001 Integration:} Reverted to single-tenant architecture. Removed organization isolation logic to simplify offline sync and deployment for remote sites. \\
    \bottomrule
\end{tabularx}

% -------------------------------------------------------
% TABLE OF CONTENTS
% -------------------------------------------------------
\newpage
\tableofcontents
\newpage

% -------------------------------------------------------
% 1. INTRODUCTION
% -------------------------------------------------------
\section{Introduction}

\subsection{Purpose}
This Software Requirements Specification (SRS) describes the complete
functional and non-functional requirements for \textbf{Beacon}, an
open-source, self-hostable, location-based attendance tracking system
designed for environments with sparse or no internet connectivity.
The system comprises an Android mobile application for field workers and
a web-based dashboard for managers and administrators.

This document serves as the single source of truth for all stakeholders,
including those involved in design, development, testing, deployment, and
community contribution. It is published openly alongside the source code
so that contributors, deployers, and evaluators can understand the
system's intent and constraints without requiring access to any private
materials.

Beacon is designed to be deployed by any organisation that needs
offline-capable, GPS-verified attendance tracking. While the reference
deployment targets a remote farm, the system is intentionally
locale-agnostic and configurable: default values reflect the reference
deployment, but deployers are expected to adjust them to their
operational context. All locale-specific defaults (public holidays,
currency, timezone, language) are configurable without modifying source
code.

\subsection{Document Conventions}
The following conventions are used throughout this document. The key words
\textbf{SHALL}, \textbf{SHOULD}, and \textbf{MAY} are to be interpreted as
described in \href{https://datatracker.ietf.org/doc/html/rfc2119}{RFC 2119}.

\begin{description}[leftmargin=3cm, style=nextline]
    \item[\textbf{FR-XX}]   Functional Requirement identifier (e.g., FR-01, FR-02).
    \item[\textbf{NFR-XX}]  Non-Functional Requirement identifier.
    \item[\textbf{BR-XX}]   Business Rule identifier.
    \item[\textbf{Bold}]    Used for requirement identifiers, field names, and key terms.
    \item[\textit{Italic}]  Used for priority labels and notes.
    \item[\texttt{Monospace}] Used for technical identifiers, API fields, and code references.
    \item[SHALL]            Denotes a mandatory requirement.
    \item[SHOULD]           Denotes a recommended but optional requirement.
    \item[MAY]              Denotes a purely optional feature.
\end{description}

Requirements are assigned one of three priority levels:
\begin{itemize}
    \item \textbf{High} --- Must be implemented in Version 1.0.
    \item \textbf{Medium} --- Strongly recommended; may be deferred to a
          subsequent sprint.
    \item \textbf{Low} --- Optional enhancement, targeted for a future
          community milestone.
\end{itemize}

\newpage

\subsection{Intended Audience and Reading Suggestions}
\begin{table}[h!]
    \centering
    \caption{Intended Audience}
    \label{tab:audience}
    \begin{tabularx}{\linewidth}{@{} l X @{}}
        \toprule
        \textbf{Audience}           & \textbf{Recommended Sections} \\
        \midrule
        Deploying Organisation      & \S1, \S2, \S3 (feature descriptions only), \S5 \\
        Software Developers         & All sections, especially \S3 and \S4 \\
        QA / Testers                & \S3 (Acceptance Criteria), \S5, Appendix B \\
        System Administrators       & \S2.4, \S4, \S5.3, \S6.2 \\
        Open-Source Contributors    & \S1.4 (Contribution Roadmap), \S3, \S4,
                                      Appendix B (Issues List) \\
        \bottomrule
    \end{tabularx}
\end{table}

\subsection{Project Scope}
\label{sec:scope}

Beacon is a two-component management system built for environments where
internet connectivity is unreliable or intermittent. It is designed to be
self-hosted: a single \texttt{docker compose up} from the repository root
brings up a fully functional instance.

\subsubsection*{Deployment Model}
\begin{description}[leftmargin=3.5cm, style=nextline]
    \item[Single-Tenant (Self-Hosted)]
        The system is deployed as a standalone instance. All users, data, and configurations belong to a single environment. This allows for simplified offline synchronization and local-network operation in remote areas without requiring a persistent connection to a central multi-tenant cloud.
\end{description}

\subsubsection*{In Scope for Version 1.0}
\begin{itemize}
    \item GPS geofence based clock-in and clock-out for workers through
          their installed Android app.
    \item Selfie photograph capture at clock-in for identity verification.
    \item Offline-first operation with background record-synchronisation
          upon restoring network access.
    \item Random mid-day GPS presence checks (frequency remotely
          configurable) during active work sessions.
    \item Worker-facing attendance history view in the Android app.
    \item Pluggable wage calculation engine with three built-in strategies:
          daily fixed rate, monthly salary divided by working days, and
          hourly rate. Deployers may register custom strategies.
    \item Remote configuration of all operational parameters: acceptable
          work locations (geofence), required GPS accuracy threshold,
          required daily work duration, and random location check frequency.
    \item Holiday management: a locale-configurable default holiday list
          (India gazetted holidays shipped as the reference preset);
          per-employee configurable weekly day(s) off; admin ability to
          mark an absence as paid or unpaid leave.
    \item Petty cash ledger: admin-controlled balance per worker; worker
          expense entry; real-time balance and transaction history visible
          in the Android app; all admin actions fully audited.
    \item Pluggable push notification provider (Firebase Cloud Messaging
          as the default implementation; alternative providers
          such as ntfy or Gotify can be substituted via configuration).
    \item Optional crash reporting and diagnostics provider (Firebase
          Crashlytics as the default; can be disabled or replaced without
          modifying application code).
    \item Daily clock-in reminder notification at a configurable time
          (default: 09:00 local time) on non-holiday working days.
    \item Post clock-out duration summary with a short-hours warning if
          the session was below the configured required work duration;
          corresponding flag in the admin dashboard.
    \item Web dashboard for managers and administrators: daily attendance
          view, holiday management, wage summary, petty-cash oversight,
          and report generation.
    \item Role-based access control (Field Worker, Manager, Admin).
    \item Configurable business rules: geofence tolerance, late policy,
          overtime, missed clock-out auto-session end.
\end{itemize}

\subsubsection*{Community Contribution Roadmap}

The following capabilities are out of scope for Version 1.0 but are
identified as well-defined contribution opportunities for the open-source
community. Each item is tracked as an issue in the project repository.

\begin{itemize}
    \item \textbf{Payroll integration} --- Automated salary disbursement
          or direct integration with banking or payroll APIs.
    \item \textbf{Inventory and bookkeeping} --- Material tracking,
          purchase records, or general ledger features.
    \item \textbf{In-app messaging} --- Worker-to-manager or
          worker-to-worker chat within the app.
    \item \textbf{Full HR management} --- Leave applications, performance
          reviews, document storage, or contract management.
    \item \textbf{ML-based identity verification} --- Facial recognition
          or liveness detection on selfie capture.
    \item \textbf{Statutory compliance calculations} --- Tax deductions,
          provident fund (PF), employee state insurance (ESI), or TDS
          computation. (These are jurisdiction-specific and intentionally
          left to deployer-provided plugins.)
    \item \textbf{Additional locale presets} --- Holiday calendars and
          wage rule presets for jurisdictions beyond India.
    \item \textbf{Half-day attendance marking} --- Sub-day attendance
          granularity.
    \item \textbf{Alternative notification backends} --- Ready-made
          provider implementations beyond FCM (ntfy, Gotify, APNS, etc.).
\end{itemize}

\subsection{References}
\begin{enumerate}
    \item Beacon --- Requirements Document (Draft), internal project notes,
          January 2026. \textit{[Primary source for this SRS.]}
    \item IEEE Std 830-1998, \textit{IEEE Recommended Practice for Software
          Requirements Specifications}. IEEE, 1998.
    \item Android Developer Documentation: Background Location Access.
          \url{https://developer.android.com/training/location}
    \item Android Developer Documentation: WorkManager.
          \url{https://developer.android.com/topic/libraries/architecture/workmanager}
    \item OWASP Mobile Security Testing Guide (MSTG), v1.7.0.
          \url{https://owasp.org/www-project-mobile-security-testing-guide/}
    \item RFC 7519: JSON Web Token (JWT). IETF, 2015.
    \item Ministry of Personnel, Public Grievances and Pensions, Government
          of India: List of Gazetted Public Holidays (annual circular).
          \textit{[Shipped as the reference locale preset; not a hard
          dependency of the system.]}
    \item Firebase Cloud Messaging (FCM) Documentation.
          \url{https://firebase.google.com/docs/cloud-messaging}
          \textit{[Default notification provider; optional.]}
\end{enumerate}

% -------------------------------------------------------
% 2. OVERALL DESCRIPTION
% -------------------------------------------------------
\section{Overall Description}

\subsection{Product Perspective}
Beacon is a standalone, purpose-built system and is not a module or
extension of any existing platform.

The system consists of two interacting components:

\begin{enumerate}
    \item \textbf{Android Worker App} --- A lightweight native Android
          application installed on field workers' devices. It handles
          authentication, GPS capture, selfie recording, and local storage
          of attendance events, syncing to the backend when connectivity is
          available.
    \item \textbf{Web Dashboard (Manager / Admin)} --- A browser-based
          interface for management staff to monitor attendance in real time,
          configure schedules and policies, manage staff accounts, and
          export reports.
\end{enumerate}

Both components communicate with a \textbf{central backend server} via a
REST API over HTTPS. The Android app is designed to function independently
of this connection for all core attendance actions, queuing records locally
until sync is possible.

\subsubsection*{Locale and Internationalisation}

Beacon ships with a reference configuration targeting Indian farm
deployments (IST timezone, INR currency, India gazetted holidays, English
UI). All locale-sensitive values are configurable without modifying source
code:

\begin{itemize}
    \item \textbf{Timezone} --- configurable per organisation; all
          internal timestamps are stored in UTC.
    \item \textbf{Currency} --- configurable; used for display in wage
          summaries and petty-cash screens.
    \item \textbf{Holiday presets} --- the India gazetted holiday list
          is the default seed; administrators may replace or extend it.
          Additional locale presets may be contributed by the community
          (see \S1.4 Contribution Roadmap).
    \item \textbf{Reminder notification time} --- default 09:00 local
          time; configurable per organisation.
\end{itemize}

No assumption is made that deployers are located in India. The India
defaults are a convenience for the reference deployment, not a
requirement of the system.

\subsection{Product Features}
A high-level summary of key features is listed below. Detailed
specifications appear in Section~\ref{sec:features}.

\begin{itemize}
    \item \textbf{Geofenced Clock-in / Clock-out} with GPS accuracy
          validation.
    \item \textbf{Selfie Capture} for identity assurance at clock-in.
    \item \textbf{Offline-First Operation} with reliable queued sync.
    \item \textbf{Random Mid-Day GPS Checks} (remotely configurable
          frequency) for continued presence verification.
    \item \textbf{Duplicate Attendance Prevention} enforced on device
          and server.
    \item \textbf{Worker Attendance History} viewable on the Android app.
    \item \textbf{Pluggable Wage Calculation Engine} with built-in
          strategies: daily fixed, monthly divided, and hourly.
    \item \textbf{Remote System Configuration} of geofence locations,
          GPS accuracy threshold, required work duration, and ping
          frequency.
    \item \textbf{Holiday Management}: locale-configurable default
          holiday list; per-worker weekly off; absentia marking as
          paid/unpaid.
    \item \textbf{Petty-Cash Ledger} with worker expense entry, balance
          view, admin top-up, and full audit trail.
    \item \textbf{Pluggable Push Notifications} (FCM default; swappable).
    \item \textbf{Optional Crash Reporting} (Crashlytics default;
          disable-able).
    \item \textbf{Short-Hours Warning} on clock-out and dashboard flag.
    \item \textbf{Manager Dashboard} with daily view, history, anomaly
          flagging, and under-hours indicators.
    \item \textbf{Staff Scheduling} and \textbf{Holiday Management} via
          calendar interface.
    \item \textbf{Report Export} in CSV and PDF formats.
    \item \textbf{Attendance Editing and Overrides} with full audit trail.
    \item \textbf{Configurable Business Rules} for geofence, late policy,
          and overtime.
\end{itemize}

\newpage
\subsection{User Classes and Characteristics}
\begin{table}[h!]
    \centering
    \caption{User Classes and Characteristics}
    \label{tab:users}
    \begin{tabularx}{\linewidth}{@{} l l X @{}}
        \toprule
        \textbf{User Class} & \textbf{Platform} & \textbf{Characteristics} \\
        \midrule
        Field Worker  & Android App      & Field staff with basic smartphone
                                           literacy. May have limited proficiency
                                           in the deployment language. Primary
                                           concern is a simple, fast clock-in/out
                                           flow. \\
        \midrule
        Manager       & Web Dashboard    & Supervisory staff who monitor daily
                                           attendance, approve anomalies, and
                                           generate reports. Moderate technical
                                           proficiency. \\
        \midrule
        Administrator & Web Dashboard    & System configurator. Sets geofence
                                           parameters, manages staff accounts and
                                           schedules. High technical proficiency
                                           expected. May be the same person as
                                           the Manager. \\
        \midrule
        Deployer      & Server / CLI     & Technical operator responsible for
                                           self-hosting the system. Expected to
                                           be comfortable with Docker and basic
                                           Linux administration. Interacts with
                                           the system via \texttt{docker compose},
                                           environment variables, and the admin
                                           dashboard. \\
        \bottomrule
    \end{tabularx}
\end{table}

\subsection{Operating Environment}
\begin{description}[leftmargin=3.5cm, style=nextline]
    \item[Android App]
        Android 10 (API level 29) or higher. Minimum 2~GB RAM. Requires
        camera, GPS (fine location), and background location permissions.
        Must operate in GPS-limited rural environments; no assumption of
        stable data connectivity. Google Play Services are required only
        when Firebase is used as the notification provider; deployments
        using an alternative provider do not have this dependency.

    \item[Web Dashboard]
        Modern evergreen browsers: Google Chrome 110+, Mozilla Firefox 110+,
        Microsoft Edge 110+. Minimum screen resolution 1280 $\times$ 720.
        Stable broadband or mobile internet connection recommended.

    \item[Backend Server]
        Linux-based server environment (Ubuntu 22.04 LTS or equivalent).
        Accessible via HTTPS (TLS 1.2 or higher). Must support concurrent
        connections from all enrolled devices and dashboard users.
        Distributed as a Docker image; no other runtime dependencies are
        required on the host. For remote sites, the server MAY be hosted on a local network (Intranet) to allow workers to sync data without external internet access.

    \item[Network]
        The Android app SHALL operate fully offline. Sync is triggered
        automatically when any network interface (WiFi or mobile data)
        becomes available. No minimum bandwidth is assumed for field
        devices.
\end{description}

\subsection{Design and Implementation Constraints}
\begin{enumerate}
    \item The Android application SHALL be developed natively for Android
          (Kotlin); cross-platform frameworks are discouraged due to GPS
          and background task reliability requirements.
    \item All data transmission between client and server SHALL use HTTPS
          with TLS 1.2 or higher.
    \item Selfie images SHALL be compressed before transmission to minimise
          bandwidth consumption (target: under 200~KB per image).
    \item The system SHALL NOT transmit raw GPS coordinates to any
          third-party analytics service.
    \item Background GPS ping frequency SHALL be configurable and SHALL
          NOT exceed 4 pings per active work session to preserve battery
          life.
    \item All attendance records SHALL store timestamps in UTC; display
          conversion to local time is the responsibility of the client
          layer.
    \item The web dashboard SHALL be accessible without installing any
          browser plugin or extension.
    \item Push notification and crash reporting integrations SHALL be
          implemented behind provider interfaces. The default
          implementations use Firebase; these MAY be replaced by
          alternative implementations via environment variable
          configuration. The core application SHALL NOT contain
          unconditional references to any specific provider SDK.
    \item The wage calculation engine SHALL be implemented as a set of
          named, selectable strategies. Built-in strategies are described
          in FR-13. Custom strategies MAY be registered by deployers
          without modifying core application code.
\end{enumerate}

\subsection{User Documentation}
The following documentation SHALL be delivered alongside the system:

\begin{itemize}
    \item \textbf{Field Worker Quick-Start Card}: A single-page visual
          guide (printable) covering app installation, login, clock-in,
          and clock-out.
    \item \textbf{Manager / Admin User Manual}: A full guide covering
          all dashboard functions: attendance monitoring, scheduling,
          holiday management, exports, and system configuration.
    \item \textbf{Self-Hosting Guide}: Covers Docker-based deployment,
          environment variable reference, backup procedures, and
          troubleshooting. Targeted at deployers.
    \item \textbf{Contributor Guide}: Covers repository structure,
          development environment setup, architecture decision records,
          and contribution workflow. Published in the repository root as
          \texttt{CONTRIBUTING.md}.
    \item \textbf{In-App Help}: Contextual error messages on the Android
          app SHALL be clear and actionable (see FR-02 acceptance
          criteria).
\end{itemize}

\subsection{Assumptions and Dependencies}

\textbf{Assumptions:}
\begin{enumerate}
    \item Each field worker is issued or personally owns an Android device
          meeting the minimum hardware specification.
    \item The work premises have a definable GPS boundary suitable for
          geofence configuration.
    \item The deploying organisation will provide accurate GPS coordinates
          for all acceptable work locations prior to system configuration.
    \item Workers will carry their assigned device during all working
          hours.
    \item At least one manager or administrator will have reliable internet
          access to operate the web dashboard.
    \item The deploying organisation will select an appropriate wage
          calculation strategy and configure it before the system goes
          live.
    \item The default holiday list (India gazetted holidays) will be
          reviewed and adjusted by the administrator to match the
          deployment locale before going live.
    \item Petty-cash balances are managed exclusively by the admin;
          workers can only add expense records, not modify balances
          directly.
\end{enumerate}

\textbf{Dependencies:}
\begin{enumerate}
    \item When Firebase is used as the notification provider, Google Play
          Services must be available on all worker devices. This dependency
          is lifted when an alternative notification provider is configured.
    \item The backend server must be deployed and reachable over the
          internet before the Android app can perform any sync operation.
    \item A push notification infrastructure endpoint (FCM or equivalent)
          is required for admin notifications, and optionally worker notifications in the future.
    \item An accurate system clock (NTP-synchronised) is assumed on the
          backend server for all timestamp generation.
\end{enumerate}

% -------------------------------------------------------
% 3. SYSTEM FEATURES
% -------------------------------------------------------
\section{System Features}
\label{sec:features}

\frblock{FR-01: Employee Authentication}{}

\srsfield{Description}
The system SHALL allow a field worker to log in using a username and
password assigned by the administrator. Each device SHALL be associated
with exactly one employee account, and the device ID and model SHALL be
recorded for audit purposes.

\srsfield{Acceptance Criteria}
\begin{itemize}
    \item Login succeeds only for accounts with valid, active credentials.
    \item On successful login, the device ID is recorded server-side.
    \item The device shall remain logged in even when it's offline for a set number of days, which SHALL be defined in the configuration, following which it has to re-authenticate with the server.
\end{itemize}

% ----------------------------------------------------------------
\subsection{Feature 2: Clock-In and Clock-Out}
% ----------------------------------------------------------------

\subsubsection*{Description and Priority}
Core attendance recording functionality. Workers record the start and end
of their work session via the Android app with GPS and photo verification.
\hfill \textit{Priority: High}

\subsubsection*{Functional Requirements}

\frblock{FR-02: Clock-In (Attendance Start)}{}

\srsfield{Description}
The system SHALL allow a logged-in field worker to initiate a clock-in
event for the current working day, subject to geofence and GPS accuracy
validation.

\srsfield{Rules}
\begin{itemize}
    \item The worker SHALL be within the configured work premises geofence
          boundary (within the allowed tolerance).
    \item A selfie photograph SHALL be captured at the moment of clock-in.
    \item Clock-in SHALL function without an active internet connection.
\end{itemize}

\srsfield{Captured Data}
\begin{itemize}
    \item \texttt{timestamp} (UTC)
    \item \texttt{gps\_lat}, \texttt{gps\_lon} (decimal degrees)
    \item \texttt{gps\_accuracy} (metres)
    \item \texttt{selfie\_image\_url} (populated after sync)
    \item \texttt{device\_id}, \texttt{app\_version}
\end{itemize}

\srsfield{Acceptance Criteria}
\begin{itemize}
    \item Clock-in succeeds only if the worker is inside the geofence
          AND GPS accuracy is within the configured threshold.
    \item If offline, the record is stored in local device storage and
          automatically synced upon network restoration.
    \item On failure, the app SHALL display the specific reason
          (e.g., ``Outside geofence'', ``GPS accuracy too low'',
          ``Camera permission denied'').
\end{itemize}

\bigskip
\frblock{FR-03: Clock-Out (Attendance End)}{}

\srsfield{Description}
The system SHALL allow a worker with an active clock-in record to record
the end of their work session.

\srsfield{Rules}
\begin{itemize}
    \item A clock-out SHALL only be possible if a clock-in exists for
          the current day.
    \item GPS coordinates SHALL be captured at clock-out.
    \item Selfie capture at clock-out is configurable (see open issues,
          Appendix~\ref{app:issues}).
    \item Clock-out location enforcement (geofence) is configurable.
\end{itemize}

\srsfield{Acceptance Criteria}
\begin{itemize}
    \item The clock-out button SHALL only be visible when an active
          clock-in record exists for the day.
    \item Offline clock-out is stored locally and synced upon
          network restoration.
\end{itemize}

% ----------------------------------------------------------------
\subsection{Feature 3: Duplicate Attendance Prevention}
% ----------------------------------------------------------------

\subsubsection*{Description and Priority}
The system prevents workers from clocking in or out more than once per
day without managerial override.
\hfill \textit{Priority: High}

\subsubsection*{Functional Requirements}

\frblock{FR-04: Prevent Duplicate Attendance}{}

\srsfield{Acceptance Criteria}
\begin{itemize}
    \item Once a clock-in record exists for a given date, further
          clock-ins on that date SHALL be blocked unless a manager
          explicitly overrides.
    \item Clock-out SHALL be permitted only once per clock-in record.
    \item Duplicate prevention SHALL be enforced both on-device and
          server-side (to handle sync edge cases).
\end{itemize}

% ----------------------------------------------------------------
\subsection{Feature 4: Random Mid-Day GPS Checks}
% ----------------------------------------------------------------

\subsubsection*{Description and Priority}
The system performs randomized background GPS pings during an active work
session to provide continued presence assurance without requiring manual
action from the worker.
\hfill \textit{Priority: Medium}

\subsubsection*{Functional Requirements}

\frblock{FR-05: Random Mid-Day GPS Checks}{}

\srsfield{Rules}
\begin{itemize}
    \item Pings SHALL only occur during an active work session
          (between clock-in and clock-out).
    \item Ping frequency SHALL be configurable (e.g., 1--4 pings per
          day).
    \item Each ping records: GPS coordinates, GPS accuracy, and timestamp.
    \item Ping behaviour SHALL be designed to minimise battery drain.
\end{itemize}

\srsfield{Acceptance Criteria}
\begin{itemize}
    \item GPS pings are queued offline and synced upon connectivity.
    \item If background location permissions or OS restrictions prevent
          pings, the system SHALL log the reason server-side.
\end{itemize}

% ----------------------------------------------------------------
\subsection{Feature 5: Manager Dashboard --- Daily Attendance View}
% ----------------------------------------------------------------

\subsubsection*{Description and Priority}
Managers can view a consolidated daily attendance summary for all staff,
including flags for anomalies.
\hfill \textit{Priority: High}

\subsubsection*{Functional Requirements}

\frblock{FR-06: Daily Attendance View}{}

\srsfield{Description}
The web dashboard SHALL present a daily attendance list showing, for each
employee: clock-in time, clock-out time, total hours worked, and any
flags or issues (e.g., outside geofence, low GPS accuracy, missed
clock-out).

\srsfield{Acceptance Criteria}
\begin{itemize}
    \item The attendance list SHALL load within 2 seconds for up to
          500 employee records on a standard broadband connection.
    \item The list SHALL be filterable by date, employee name, and
          flag status.
\end{itemize}

% ----------------------------------------------------------------
\subsection{Feature 6: Individual Staff Attendance History}
% ----------------------------------------------------------------

\subsubsection*{Description and Priority}
Managers can drill down into a per-employee attendance history view.
\hfill \textit{Priority: High}

\subsubsection*{Functional Requirements}

\frblock{FR-07: Staff Attendance History}{}

\srsfield{Description}
The dashboard SHALL allow managers to view a selected employee's
attendance history, including daily records, total hours, anomaly flags,
and selfie photo references (subject to privacy configuration).

% ----------------------------------------------------------------
\subsection{Feature 7: Staff Scheduling}
% ----------------------------------------------------------------

\subsubsection*{Description and Priority}
Administrators and managers can configure working day patterns per
employee.
\hfill \textit{Priority: Medium}

\subsubsection*{Functional Requirements}

\frblock{FR-08: Staff Scheduling (Working Days per Employee)}{}

\srsfield{Rules}
\begin{itemize}
    \item Each employee SHALL have a configurable working day pattern
          (Monday--Sunday toggles or a shift template).
    \item Schedule changes SHALL apply to future dates by default;
          retroactive application is configurable.
\end{itemize}

\srsfield{Acceptance Criteria}
\begin{itemize}
    \item Schedule changes are immediately reflected in the attendance
          dashboard view.
    \item If a worker clocks in on a non-working day, the record is
          flagged automatically.
\end{itemize}

% ----------------------------------------------------------------
\subsection{Feature 8: Holiday Management}
% ----------------------------------------------------------------

\subsubsection*{Description and Priority}
Administrators can define public and organisation-specific holidays via a
calendar interface.
\hfill \textit{Priority: Medium}

\subsubsection*{Functional Requirements}

\frblock{FR-09: Public Holiday Management}{}

\srsfield{Rules}
\begin{itemize}
    \item Holidays SHALL be markable on a calendar interface.
    \item Each holiday SHALL be configurable as paid or unpaid.
    \item Holidays SHALL be applicable to all employees or a selected
          subset.
\end{itemize}

\srsfield{Acceptance Criteria}
\begin{itemize}
    \item Holidays are visible in the dashboard calendar view.
    \item Attendance recorded on a holiday is flagged or handled
          according to the configured rule.
\end{itemize}

% ----------------------------------------------------------------
\subsection{Feature 9: Reports and Export}
% ----------------------------------------------------------------

\subsubsection*{Description and Priority}
Managers can generate and export attendance reports for payroll and
compliance purposes.
\hfill \textit{Priority: High}

\subsubsection*{Functional Requirements}

\frblock{FR-10: Attendance Report Export}{}

\srsfield{Rules}
\begin{itemize}
    \item Reports SHALL be exportable in CSV and PDF formats.
    \item Export SHALL be filterable by date range and/or individual
          employee.
    \item Reports SHALL include: days present, days absent, total hours,
          and overtime hours.
\end{itemize}

% ----------------------------------------------------------------
\subsection{Feature 10: Attendance Edits and Overrides}
% ----------------------------------------------------------------

\subsubsection*{Description and Priority}
Managers or administrators can correct erroneous records or add missing
entries, with a complete audit trail.
\hfill \textit{Priority: Medium}

\subsubsection*{Functional Requirements}

\frblock{FR-11: Attendance Edits and Overrides}{}

\srsfield{Rules}
\begin{itemize}
    \item Authorised users SHALL be able to add missing entries, correct
          existing entries, and approve flagged anomalies.
    \item Every edit SHALL store: editor identity, edit timestamp, and a
          mandatory reason note.
    \item The original record SHALL be preserved for audit; edits create
          a versioned history.
\end{itemize}

% ----------------------------------------------------------------
\subsection{Feature 11: Worker Attendance History (Mobile)}
% ----------------------------------------------------------------

\subsubsection*{Description and Priority}
Field workers can view their own attendance history directly within the
Android app, without requiring access to the web dashboard.
\hfill \textit{Priority: High}

\subsubsection*{Functional Requirements}

\frblock{FR-12: Worker Self-Service Attendance History}{}

\srsfield{Description}
The Android app SHALL provide a dedicated screen showing the authenticated
worker's attendance records, including date, clock-in time, clock-out
time, total hours worked, holiday status, and any flags (e.g., short
hours, outside geofence).

\srsfield{Rules}
\begin{itemize}
    \item History SHALL be browsable by month with a summary of days
          present, days absent, and total hours.
    \item Records SHALL be read-only for the worker; no editing is
          permitted from this view.
    \item History SHALL be available offline for previously synced records.
\end{itemize}

\srsfield{Acceptance Criteria}
\begin{itemize}
    \item The attendance history screen loads within 2 seconds for up to
          90 days of records from local cache.
    \item Workers can only view their own records; cross-worker data
          access is not permitted.
\end{itemize}

% ----------------------------------------------------------------
\subsection{Feature 12: Wage Calculation Engine}
% ----------------------------------------------------------------

\subsubsection*{Description and Priority}
The system automatically calculates each worker's earned wages for a
given period using a configurable, pluggable calculation strategy.
\hfill \textit{Priority: High}

\subsubsection*{Functional Requirements}

\frblock{FR-13: Pluggable Wage Calculation Engine}{}

\srsfield{Description}
The system SHALL calculate a worker's gross wages for any date range
using the strategy configured for that worker. Each worker SHALL have a
\texttt{wage\_strategy} field that selects which calculation strategy to
apply. Three strategies are built in; deployers MAY register additional
strategies without modifying core application code.

\srsfield{Built-in Wage Strategies}

\begin{description}[leftmargin=4cm, style=nextline]
    \item[\texttt{DAILY\_FIXED}]
        \textit{(Default)} The worker has a configured daily rate.
        Wage = paid days $\times$ daily rate. A day is a ``paid day'' if
        the worker attended, or if the day is a paid public holiday, or
        if the admin has marked the absence as paid leave. Weekly off
        days and unpaid absences are excluded.

    \item[\texttt{MONTHLY\_DIVIDED}]
        The worker has a configured monthly salary. For a given month,
        the effective daily rate = monthly salary $\div$ number of
        calendar days in that month. Wage = effective daily rate
        $\times$ paid days. Useful when the employment contract specifies
        a fixed monthly figure rather than a per-day rate.

    \item[\texttt{HOURLY}]
        The worker has a configured hourly rate. Wage = hourly rate
        $\times$ total hours worked, derived from session durations in
        the attendance records. Overtime hours (if the overtime policy
        is enabled) MAY be calculated at a separate configured multiplier.
\end{description}

\srsfield{Rules (all strategies)}
\begin{itemize}
    \item Weekly off days are excluded from all wage calculations
          regardless of strategy.
    \item Wage summaries SHALL be viewable by the worker on the Android
          app and by managers/admins on the dashboard.
    \item The wage calculation SHALL be performed server-side;
          client-side figures are display-only and refreshed on each sync.
\end{itemize}

\srsfield{Acceptance Criteria}
\begin{itemize}
    \item Wage calculations update automatically when attendance records,
          leave markings, or rate/salary values are modified.
    \item The Android app displays the worker's current month wage
          summary on the home screen or a dedicated earnings screen.
    \item Managers can export wage summaries per employee or for all
          employees for a selected date range (CSV / PDF).
    \item Switching a worker's strategy recalculates wages from the
          effective date of the change; historical records are not
          retroactively recalculated unless explicitly triggered by
          an admin.
\end{itemize}

% ----------------------------------------------------------------
\subsection{Feature 13: Remote System Configuration}
% ----------------------------------------------------------------

\subsubsection*{Description and Priority}
Administrators can remotely configure all operational parameters of the
system from the web dashboard. Changes propagate to worker devices on
next sync without requiring an app update.
\hfill \textit{Priority: High}

\subsubsection*{Functional Requirements}

\frblock{FR-14: Remote Configuration Panel}{}

\srsfield{Description}
The admin dashboard SHALL provide a configuration panel through which
all system-wide and per-worker operational parameters can be updated
remotely.

\srsfield{Configurable Parameters}
\begin{itemize}
    \item \textbf{Acceptable Work Locations:} One or more named geofence
          zones (radius or polygon), each with a centre-point, radius,
          and boundary tolerance. Multiple zones may be active
          simultaneously.
    \item \textbf{Required GPS Accuracy Threshold:} Maximum acceptable
          GPS error radius (metres) for clock-in acceptance.
    \item \textbf{Required Daily Work Duration:} Minimum hours per day
          that constitute a full attendance; used for short-hours
          detection.
    \item \textbf{Random Location Check Frequency:} Number of background
          GPS pings per active work session (range: 0--4).
    \item \textbf{Offline Login Window:} Maximum duration (days) for
          cached credential login without server re-authentication.
    \item \textbf{Missed Clock-Out Auto-Close:} Enable/disable and
          threshold in hours.
    \item \textbf{Late Clock-In Policy:} Threshold time and grace period.
    \item \textbf{Overtime Policy:} Daily threshold and rounding rule.
    \item \textbf{Wage Strategy (per worker):} One of \texttt{DAILY\_FIXED},
          \texttt{MONTHLY\_DIVIDED}, \texttt{HOURLY}, or a custom
          registered strategy identifier.
    \item \textbf{Notification Provider:} Identifier of the active push
          notification provider (e.g., \texttt{fcm}, \texttt{ntfy}).
          Read-only on the dashboard; set via server environment
          configuration.
\end{itemize}

\srsfield{Rules}
\begin{itemize}
    \item Configuration changes SHALL be versioned and timestamped.
    \item The Android app SHALL fetch the latest configuration on every
          successful sync and apply it immediately.
    \item Configuration is applied system-wide unless a parameter
          supports per-worker override (noted per parameter above).
    \item All configuration changes SHALL be recorded in the audit log
          with the admin identity.
\end{itemize}

\srsfield{Acceptance Criteria}
\begin{itemize}
    \item A configuration change made on the dashboard is reflected on
          a connected worker device within one sync cycle (target:
          60 seconds on available network).
    \item Invalid configuration values (e.g., negative GPS threshold)
          are rejected with a descriptive error message.
\end{itemize}

% ----------------------------------------------------------------
\subsection{Feature 14: Holiday Management: Locale Defaults \&
Per-Worker Weekly Off}
% ----------------------------------------------------------------

\subsubsection*{Description and Priority}
The system ships with a locale-configurable default holiday list. The
reference preset contains all Government of India gazetted public
holidays and is applied by default, but administrators are expected to
review and adjust it to match their deployment context. Administrators
can additionally configure per-employee weekly off days and ad-hoc
holidays.
\hfill \textit{Priority: High}

\subsubsection*{Functional Requirements}

\frblock{FR-15: Locale-Configurable Holiday Defaults}{}

\srsfield{Description}
The system SHALL ship with a pre-seeded holiday list for the active
locale preset. The reference preset (India) contains all Government of
India gazetted national and festival public holidays for the current
calendar year, marked as paid by default.

\srsfield{Rules}
\begin{itemize}
    \item The active locale preset SHALL be configurable via the admin
          dashboard or server environment without a software release.
    \item The reference India preset SHALL include at minimum: Republic
          Day (26 Jan), Holi, Good Friday, Ambedkar Jayanti (14 Apr),
          Labour Day (1 May), Independence Day (15 Aug), Gandhi Jayanti
          (2 Oct), Dussehra, Diwali, and Christmas Day (25 Dec).
    \item Administrators SHALL be able to add, remove, or reclassify any
          holiday (paid $\leftrightarrow$ unpaid) for their organisation.
    \item Holiday applicability can be set to: all employees, or a named
          subset.
    \item Additional locale presets (holiday lists for other
          jurisdictions) MAY be contributed to the repository and
          selected via configuration.
\end{itemize}

\srsfield{Acceptance Criteria}
\begin{itemize}
    \item On a configured paid holiday, a worker who does not clock in
          SHALL NOT be flagged as absent; the day is recorded as a paid
          holiday automatically.
    \item The holiday calendar is visible to workers on the Android app
          (read-only).
\end{itemize}

\bigskip
\frblock{FR-16: Per-Worker Weekly Off Configuration}{}

\srsfield{Description}
Each worker's weekly day(s) off SHALL be individually configurable by
the administrator. Weekly off days are excluded from attendance
expectations and wage calculations.

\srsfield{Rules}
\begin{itemize}
    \item Default weekly off: Sunday (configurable to any day or
          combination of days).
    \item Changes to a worker's weekly off pattern apply from the
          effective date forward; historical records are not retroactively
          modified.
    \item If a public holiday falls on a weekly off day, the holiday
          takes precedence in display; no additional compensation is
          implied by this SRS.
\end{itemize}

\srsfield{Acceptance Criteria}
\begin{itemize}
    \item A worker's weekly off day does not trigger an absent flag or
          a clock-in reminder notification.
    \item Weekly off days are visually distinguished in the worker's
          attendance history view.
\end{itemize}

\bigskip
\frblock{FR-17: Absentia Marking as Paid / Unpaid Leave}{}

\srsfield{Description}
When a worker is absent on a working day (not a public holiday or weekly
off), the administrator SHALL be able to retroactively mark the absence
as either paid leave or unpaid leave.

\srsfield{Rules}
\begin{itemize}
    \item Absent days with no marking default to \textit{unpaid}.
    \item A paid leave marking causes the day to count as a paid day
          in wage calculations.
    \item An unpaid leave marking explicitly records the day as excluded
          from wages.
    \item All leave markings SHALL be recorded in the audit log with
          admin identity, timestamp, and optional note.
    \item The worker's attendance history view in the Android app SHALL
          reflect the leave classification (paid leave / unpaid leave /
          absent) once synced.
\end{itemize}

\srsfield{Acceptance Criteria}
\begin{itemize}
    \item Leave marking updates wage calculations immediately upon save.
    \item The worker can see the leave classification in their attendance
          history; they cannot modify it.
\end{itemize}

% ----------------------------------------------------------------
\subsection{Feature 15: Petty-Cash Ledger}
% ----------------------------------------------------------------

\subsubsection*{Description and Priority}
The system maintains a per-worker petty-cash balance. Workers can record
individual expenditures; administrators control the balance and can view
all transactions. All admin actions are audited and visible as entries in
the worker's transaction feed.
\hfill \textit{Priority: High}

\subsubsection*{Functional Requirements}

\frblock{FR-18: Worker Petty-Cash View and Expense Entry}{}

\srsfield{Description}
The Android app SHALL display the worker's current petty-cash balance and
a chronological list of all transactions (credits and debits), including
admin-initiated entries.

\srsfield{Rules}
\begin{itemize}
    \item Workers can create \textit{expense records} (debit entries)
          containing: amount, date, and a short description.
    \item Workers can view their current balance and full transaction
          history.
    \item Workers CANNOT modify or delete existing entries.
    \item Admin-initiated entries (top-ups, adjustments, corrections)
          SHALL appear in the worker's transaction feed with a system
          label (e.g., ``Admin credit --- [note]'').
    \item The petty-cash screen SHALL function offline; expense records
          are queued and synced.
\end{itemize}

\srsfield{Captured Data per Expense Entry}
\begin{itemize}
    \item \texttt{entry\_id}, \texttt{worker\_id}
    \item \texttt{entry\_type}: \texttt{debit} (worker expense) or
          \texttt{credit} (admin top-up / adjustment)
    \item \texttt{amount} (positive decimal)
    \item \texttt{description} (free text, max 200 characters)
    \item \texttt{entry\_date} (worker-entered date; defaults to today)
    \item \texttt{created\_at} (UTC timestamp of record creation)
    \item \texttt{created\_by}: worker ID or admin ID
    \item \texttt{audit\_note} (mandatory for admin-initiated entries)
\end{itemize}

\srsfield{Acceptance Criteria}
\begin{itemize}
    \item The balance displayed is the sum of all credits minus all
          debits for the worker.
    \item New worker expense entries appear immediately in the local
          transaction list; they are marked as ``pending sync'' until
          server acknowledgement.
    \item Admin entries appear in the worker's feed within one sync
          cycle of the admin action.
\end{itemize}

\bigskip
\frblock{FR-19: Admin Petty-Cash Management}{}

\srsfield{Description}
Administrators SHALL be able to view all workers' petty-cash balances
and transaction histories, and to update any worker's balance via credit
or adjustment entries, each with a mandatory audit note.

\srsfield{Rules}
\begin{itemize}
    \item The dashboard SHALL show a petty-cash summary table: all
          workers, current balance, last transaction date.
    \item Admin can drill into any worker's full transaction history.
    \item Admin can add a credit entry (top-up) or a corrective
          adjustment (positive or negative) with a mandatory note.
    \item Admin CANNOT delete worker expense entries; corrections must
          be made via adjustment entries.
    \item Every admin action (top-up, adjustment, balance reset) is
          recorded in the system-wide audit log AND creates a visible
          entry in the affected worker's transaction feed.
\end{itemize}

\srsfield{Acceptance Criteria}
\begin{itemize}
    \item An admin credit entry updates the worker's displayed balance
          within one sync cycle on the worker's device.
    \item The audit log entry for every admin petty-cash action includes:
          admin identity, target worker, entry type, amount, note, and
          UTC timestamp.
    \item The petty-cash dashboard view loads within 2 seconds for up
          to 500 workers.
\end{itemize}

% ----------------------------------------------------------------
\subsection{Feature 16: Clock-In Reminder Notifications}
% ----------------------------------------------------------------

\subsubsection*{Description and Priority}
The Android app notifies workers to mark their attendance on working days
where they have not yet clocked in. Notification delivery is handled by
a pluggable provider; Firebase Cloud Messaging is the default.
\hfill \textit{Priority: Medium}

\subsubsection*{Functional Requirements}

\frblock{FR-20: Daily Clock-In Reminder}{}

\srsfield{Description}
The system SHALL send a push notification to a worker's device at the
configured reminder time (default: 09:00 local time) on any day that is
a scheduled working day (not a public holiday, not the worker's weekly
off day) if the worker has not yet clocked in for that day.

\srsfield{Rules}
\begin{itemize}
    \item Notifications SHALL be delivered via the configured push
          notification provider. The provider is selected via server
          environment configuration; the application SHALL NOT hardcode
          a provider.
    \item The notification SHALL NOT be sent if the worker has already
          clocked in for the day by the reminder time.
    \item The notification SHALL NOT be sent on public holidays or the
          worker's configured weekly off day(s).
    \item If the device is offline at the scheduled time, the
          notification SHALL be delivered as soon as the device comes
          online, provided the worker has still not clocked in for that
          day.
    \item Workers SHALL be able to disable this notification in the app
          settings (opt-out).
    \item The reminder time SHALL be configurable per organisation
          (default: 09:00 local time).
\end{itemize}

\srsfield{Acceptance Criteria}
\begin{itemize}
    \item Notification is not delivered on holidays or weekly off days.
    \item Notification is not delivered if clock-in already exists.
    \item Tapping the notification opens the clock-in screen directly.
\end{itemize}

% ----------------------------------------------------------------
\subsection{Feature 17: Short-Hours Detection and Warning}
% ----------------------------------------------------------------

\subsubsection*{Description and Priority}
Upon clock-out, the app informs the worker of their total session
duration and warns them if they have worked less than the required daily
duration. The same flag appears in the manager dashboard.
\hfill \textit{Priority: Medium}

\subsubsection*{Functional Requirements}

\frblock{FR-21: Post-Clock-Out Duration Summary and Warning}{}

\srsfield{Description}
Immediately after a successful clock-out, the Android app SHALL display
a summary screen showing the total duration of the work session. If the
duration is less than the configured required work duration (see FR-14,
BR-09), a short-hours warning SHALL be prominently displayed.

\srsfield{Rules}
\begin{itemize}
    \item Session duration = clock-out timestamp $-$ clock-in timestamp.
    \item The warning is informational only; no automated action
          (deduction, lock-out, or escalation) is triggered.
    \item The short-hours flag SHALL be recorded on the attendance record
          server-side.
    \item The manager dashboard daily attendance view SHALL display a
          visible indicator (e.g., flag or colour coding) on any record
          where the session duration was below the required threshold.
    \item The flag is visible to managers; no notification is sent to
          the manager automatically.
\end{itemize}

\srsfield{Acceptance Criteria}
\begin{itemize}
    \item The post-clock-out screen always shows total session duration
          in hours and minutes.
    \item If duration $<$ configured minimum, a warning message is shown
          before the worker can dismiss the screen.
    \item The short-hours flag is visible in the manager dashboard for
          the affected date.
    \item If the session was auto-closed (missed clock-out), the flag
          applies if the auto-close duration is less than the minimum.
\end{itemize}


% ================================================================
% SECTION 4 — EXTERNAL INTERFACE REQUIREMENTS
% ================================================================
\section{External Interface Requirements}

\subsection{User Interfaces}

\textbf{Android Worker App screens:}
\begin{itemize}
    \item \textbf{Login Screen}: Username/password fields.
    \item \textbf{Home Screen}: Today's attendance status (clocked in /
          not clocked in / clocked out), current month wage summary,
          petty-cash balance snapshot, and sync status indicator.
    \item \textbf{Clock-In Flow}: GPS acquisition indicator, selfie
          camera capture, confirmation, and error message display.
    \item \textbf{Clock-Out Flow}: GPS confirmation, optional selfie,
          post-clock-out duration summary screen with short-hours warning
          if applicable.
    \item \textbf{Attendance History Screen}: Monthly calendar / list
          view of the worker's own attendance records; read-only.
    \item \textbf{Earnings Screen}: Month-to-date wage calculation
          breakdown (days attended or hours worked depending on strategy,
          rate, total earned).
    \item \textbf{Petty-Cash Screen}: Current balance, chronological
          transaction list (worker expenses and admin entries), and
          add-expense form.
    \item \textbf{Holiday Calendar Screen}: Read-only view of
          organisation holidays and the worker's weekly off days.
    \item \textbf{Sync Status Screen}: Queue of pending records, last
          sync time, retry option.
    \item \textbf{Settings Screen}: Notification opt-out toggle and app
          information.
\end{itemize}

\textbf{Web Dashboard screens:}
\begin{itemize}
    \item \textbf{Login}: Credential entry with session management.
    \item \textbf{Daily Attendance View}: Tabular summary with filters,
          flag indicators, and short-hours markers.
    \item \textbf{Employee List}: Searchable, paginated list with account
          management actions.
    \item \textbf{Employee History}: Per-employee timeline view with
          leave classification, wage strategy, and wage summary.
    \item \textbf{Absentia / Leave Marking}: Interface to mark an absent
          day as paid or unpaid leave for a specific employee.
    \item \textbf{Wage Summary}: Per-employee and aggregate wage
          calculation view with strategy label; exportable.
    \item \textbf{Scheduling Calendar}: Visual week/month calendar with
          working day toggles per employee.
    \item \textbf{Holiday Calendar}: Date picker for marking holidays
          with type (paid/unpaid) and applicability settings; active
          locale preset holidays visible and editable.
    \item \textbf{Petty-Cash Dashboard}: Summary table of all worker
          balances; drill-down to individual transaction history; add
          credit / adjustment form.
    \item \textbf{Export / Reports}: Date range selector, format
          selector (CSV / PDF), download trigger; covers attendance,
          wages, and petty-cash.
    \item \textbf{Admin Settings / Remote Configuration}: All operational
          parameters: geofence zones, GPS accuracy threshold, required
          work duration, ping frequency, login window, policy toggles,
          wage strategies, and locale settings.
    \item \textbf{Audit Log Viewer}: Searchable, filterable log of all
          admin actions across attendance, configuration, and petty-cash.
\end{itemize}

UI design SHALL prioritise simplicity and minimal tap-count for the
worker app. Error messages SHALL be written in plain language, not
technical codes.

\subsection{Hardware Interfaces}

\begin{description}[leftmargin=3cm, style=nextline]
    \item[GPS Receiver]
        The Android app SHALL interface with the device GPS receiver via
        the Android Location API (FusedLocationProviderClient). Accuracy
        metadata SHALL be recorded with every GPS reading.
    \item[Camera]
        The app SHALL access the device camera via Android Camera2 API
        or CameraX for selfie capture. Front or rear camera selection is
        configurable (see Appendix~\ref{app:issues}, issue 5).
    \item[Storage]
        Local SQLite or Room database on the Android device SHALL be used
        to persist offline attendance records.
    \item[Server Hardware]
        No specific server hardware is mandated; the backend SHALL run on
        any Linux-based virtual machine or cloud instance meeting minimum
        CPU and memory specifications.
\end{description}

\subsection{Software Interfaces}

\begin{description}[leftmargin=3cm, style=nextline]
    \item[Android SDK]
        Minimum API level 29 (Android 10). Google Play Services required
        only when Firebase notification provider is active.
    \item[REST API]
        The Android app and web dashboard communicate with the backend
        via a RESTful API over HTTPS. Request and response bodies use
        JSON format (UTF-8 encoding).
    \item[Authentication]
        JSON Web Tokens (JWT, RFC 7519) SHALL be used for session
        management. Tokens SHALL have a configurable expiry.
    \item[Database]
        Backend database technology is not mandated by this SRS but SHALL
        support ACID-compliant transactions given the audit requirements.
    \item[Notification Provider Interface]
        The backend SHALL expose a \texttt{NotificationProvider}
        interface. The FCM implementation is bundled by default. An
        alternative implementation MAY be activated by setting the
        \texttt{BEACON\_NOTIFICATION\_PROVIDER} environment variable
        and providing the corresponding configuration.
    \item[Crash Reporting Interface]
        The Android app SHALL expose a \texttt{CrashReporter} interface.
        The Firebase Crashlytics implementation is bundled by default
        and MAY be disabled by setting
        \texttt{BEACON\_CRASH\_REPORTING=disabled}. When disabled, the
        no-op implementation is activated and no data leaves the device.
    \item[Export Libraries]
        Server-side libraries for PDF generation and CSV serialisation
        are required for FR-10.
\end{description}

\subsection{Communications Interfaces}

\begin{description}[leftmargin=3cm, style=nextline]
    \item[HTTPS / TLS]
        All client--server communication SHALL use HTTPS with TLS 1.2
        or higher. Self-signed certificates SHALL NOT be used in
        production.
    \item[Sync Protocol]
        The Android app SHALL use a background sync mechanism (WorkManager
        or equivalent) that retries failed sync attempts with exponential
        backoff, up to a maximum retry count.
    \item[Image Upload]
        Selfie images SHALL be uploaded as multipart/form-data or as
        Base64-encoded payloads within the JSON body, with a maximum
        compressed size of 200~KB per image.
    \item[Push Notifications]
        Push notifications SHALL be dispatched via the configured
        \texttt{NotificationProvider}. The default provider is Firebase
        Cloud Messaging (FCM). Provider selection does not require an
        app update; it is a server-side configuration change.
\end{description}

% -------------------------------------------------------
% 5. OTHER NON-FUNCTIONAL REQUIREMENTS
% -------------------------------------------------------
\section{Other Non-Functional Requirements}

\subsection{Performance Requirements}

\begin{description}[leftmargin=2.5cm, style=nextline]
    \item[NFR-01]
        \textbf{App Startup:} The Android app SHALL reach the home screen
        within 3 seconds on a mid-range device (2~GB RAM, Android 10,
        cold start).
    \item[NFR-02]
        \textbf{Clock-In Response:} The clock-in confirmation SHALL be
        displayed to the user within 5 seconds of tapping the clock-in
        button, assuming GPS is already acquired.
    \item[NFR-03]
        \textbf{Sync Latency:} Offline-queued records SHALL sync to the
        server within 30 seconds of network restoration on a connection
        of at least 100~kbps (2G equivalent).
    \item[NFR-04]
        \textbf{Dashboard Load:} The daily attendance view SHALL render
        within 2 seconds for up to 500 employee records on a standard
        broadband connection ($\geq$ 10~Mbps).
    \item[NFR-05]
        \textbf{Battery Impact:} Background GPS ping operations SHALL NOT
        cause measurable battery drain exceeding 5\% per hour above the
        device idle baseline.
\end{description}

\subsection{Safety Requirements}

\begin{description}[leftmargin=2.5cm, style=nextline]
    \item[NFR-06]
        \textbf{Data Loss Prevention:} No attendance record SHALL be
        permanently lost due to a sync failure. Records SHALL remain in
        local storage until server-side acknowledgement is received.
    \item[NFR-07]
        \textbf{Graceful Degradation:} If the backend server is
        unavailable, the Android app SHALL continue to accept clock-in
        and clock-out events without displaying an error to the worker
        (beyond a sync-pending indicator).
    \item[NFR-08]
        \textbf{Conflict Resolution:} In the event of a sync conflict
        (e.g., manager edited a record while an offline record exists),
        the system SHALL flag the conflict for manual review rather than
        silently overwriting either version.
\end{description}

\subsection{Security Requirements}

\begin{description}[leftmargin=2.5cm, style=nextline]
    \item[NFR-09]
        \textbf{Authentication:} All API endpoints SHALL require a valid
        JWT token. Unauthenticated requests SHALL return HTTP 401.
    \item[NFR-10]
        \textbf{Role-Based Access Control:} The system SHALL enforce role
        boundaries between Field Worker, Manager, and Admin. A field
        worker SHALL NOT be able to access any manager or admin endpoint.
    \item[NFR-11]
        \textbf{Audit Log:} All attendance edits, overrides, and
        administrative actions (account creation/deletion, config changes)
        SHALL be recorded in an immutable audit log with actor identity
        and timestamp.
    \item[NFR-12]
        \textbf{Data in Transit:} All data transmitted between any client
        and the server SHALL be encrypted using TLS 1.2 or higher.
    \item[NFR-13]
        \textbf{Data at Rest:} Selfie images and personal data stored on
        the server SHALL be stored with access controls limiting retrieval
        to authorised roles.
    \item[NFR-14]
        \textbf{Device Binding:} The system SHALL record device ID at
        login and flag anomalies where the same account logs in from a
        different device without prior de-registration.
    \item[NFR-15]
        \textbf{Petty-Cash Ledger Integrity:} The petty-cash balance
        SHALL be calculated server-side as the authoritative sum of all
        credit and debit entries. Client-side balance is for display only
        and is overwritten on each sync. No balance field SHALL be stored
        as a mutable scalar to prevent tampering.
    \item[NFR-16]
        \textbf{Wage Calculation Consistency:} Wage figures displayed on
        the worker app and manager dashboard SHALL be derived from the
        same server-side calculation engine. Discrepancies between views
        are not permitted.
    \item[NFR-17]
        \textbf{Notification Delivery:} Clock-in reminder notifications
        SHALL be delivered within 5 minutes of the scheduled reminder
        time on a device with active network connectivity.
    \item[NFR-18]
        \textbf{Configuration Propagation:} Remote configuration changes
        SHALL reach all active devices within 60 seconds of save on a
        device with network connectivity.
    \item[NFR-19]
        \textbf{Provider Isolation:} The crash reporting provider SHALL
        NOT have access to attendance records, GPS coordinates, selfie
        images, wage data, or petty-cash data. It SHALL receive only
        anonymised crash context (stack traces, device model, OS version,
        app version).
\end{description}

\subsection{Software Quality Attributes}

\begin{description}[leftmargin=3cm, style=nextline]
    \item[Reliability]
        The backend service SHALL target 99.5\% uptime (excluding
        scheduled maintenance windows). The Android app SHALL have zero
        data loss under all network conditions.
    \item[Maintainability]
        The codebase SHALL follow a documented coding standard. All
        business rule thresholds (geofence radius, ping frequency,
        late-in time, etc.) SHALL be stored as configurable parameters,
        not hard-coded values. Provider interfaces SHALL be documented
        with enough detail for a contributor to implement a new provider
        without reading core application code.
    \item[Usability]
        The full clock-in flow (app open to confirmation) SHALL be
        completable in under 60 seconds by a first-time user following
        the quick-start card.
    \item[Portability]
        The web dashboard SHALL function correctly on Chrome, Firefox,
        and Edge (current and previous major versions) without any
        browser-specific workarounds.
    \item[Debuggability]
        The Android app SHALL maintain an on-device event log (``black
        box'') exportable by administrators for support purposes. The
        backend SHALL assign correlation IDs to all requests for
        end-to-end tracing.
    \item[Scalability]
        The system architecture SHALL support scaling to at least 1,000
        concurrent active devices without architectural changes.
    \item[Self-Hostability]
        A deployer with basic Docker and Linux knowledge SHALL be able
        to bring up a fully functional instance by following the
        self-hosting guide, without writing any code or modifying any
        source files.
\end{description}

% ================================================================
% SECTION 6 — OTHER REQUIREMENTS
% ================================================================
\section{Other Requirements}

\subsection{Business Rules}

The following rules govern system behaviour. All default values are
appropriate for the reference deployment (remote farm, India) and are
expected to be reviewed and adjusted by the deploying organisation.
Every rule marked as \textit{configurable} can be changed via the admin
dashboard or server environment without modifying source code.

\begin{description}[leftmargin=3.5cm, style=nextline]
    \item[BR-01 Geofence Type]
        Geofence boundary may be defined as a radius (circular) or a
        polygon. Multiple named zones may be active simultaneously.
        Default: single radius zone. \textit{Configurable.}
    \item[BR-02 Geofence Tolerance]
        Acceptable distance outside the geofence boundary before
        clock-in is rejected. Default: 150~m. \textit{Configurable.}
    \item[BR-03 GPS Accuracy Threshold]
        Maximum acceptable GPS error radius for clock-in.
        Default: 50~m. \textit{Configurable} (FR-14).
    \item[BR-04 Clock-Out Location]
        Whether clock-out is permitted outside the geofence.
        Default: Yes (allowed anywhere). \textit{Configurable.}
    \item[BR-05 Missed Clock-Out]
        Whether to auto-close a session after a configurable number of
        hours. Default: Yes, after 10 hours. Manager notification on
        auto-close: Yes. \textit{Configurable.}
    \item[BR-06 Late Policy]
        A clock-in after a configurable time is flagged as late. Grace
        period configurable (default: 15 minutes). Optional; disabled
        by default. \textit{Configurable.}
    \item[BR-07 Overtime Policy]
        Hours worked beyond a daily threshold are flagged as overtime.
        Rounding rule configurable (e.g., nearest 15 minutes). Optional;
        disabled by default. \textit{Configurable.}
    \item[BR-08 Mid-Day Ping Frequency]
        Number of random GPS checks per active session. Range: 0--4.
        Default: 2. \textit{Configurable} (FR-14).
    \item[BR-09 Required Daily Work Duration]
        Minimum session duration constituting a full day's work. Used
        for short-hours detection (FR-21). Default: 8 hours.
        \textit{Configurable} (FR-14).
    \item[BR-10 Worker Weekly Off Day]
        Default weekly off day for all workers: Sunday. Overridable per
        worker to any day(s) of the week (FR-16). Weekly off days are
        excluded from attendance expectations and wage calculations.
        \textit{Configurable per worker.}
    \item[BR-11 Holiday Locale Preset]
        At system initialisation, the active locale preset is seeded as
        the organisation's default paid holiday list (FR-15). The
        reference preset is the India gazetted holiday list. Deployers
        in other locales should replace or extend this list before going
        live. \textit{Configurable without a software release.}
    \item[BR-12 Absent Day Default Classification]
        An absent working day with no admin leave marking defaults to
        \textit{unpaid}. Admins may reclassify to paid or unpaid leave
        (FR-17). \textit{Configurable.}
    \item[BR-13 Wage Strategy and Rate]
        Each worker has an individually configurable wage strategy
        (\texttt{DAILY\_FIXED}, \texttt{MONTHLY\_DIVIDED},
        \texttt{HOURLY}, or a custom strategy) and the corresponding
        rate or salary figure. Currency is configurable per organisation;
        default: INR. \textit{Configurable per worker.}
    \item[BR-14 Clock-In Reminder Time]
        Daily reminder notification time. Default: 09:00 local time.
        \textit{Configurable per organisation} (FR-14).
    \item[BR-15 Petty-Cash Edit Policy]
        Workers may add expense (debit) entries only. Admins may add
        credit or adjustment entries. No entry may be deleted by any
        user; corrections are made via adjustment entries (FR-19).
        \textit{Not overridable; this is a data integrity rule.}
\end{description}

\subsection{Deployment Requirements}

\begin{enumerate}
    \item The system SHALL be distributable as a set of Docker images
          orchestrated by a \texttt{docker-compose.yml} in the repository
          root. A deployer SHALL be able to bring up a fully functional
          instance with a single command.
    \item An \texttt{.env.example} file SHALL be provided listing all
          required and optional environment variables with descriptions
          and default values.
    \item The system SHALL undergo a field trial on a minimum of 1 device
          for 2--3 days before any production rollout.
    \item Full rollout SHALL proceed only after the field trial produces
          no critical defects.
    \item An installation guide SHALL be provided for sideloading the APK
          onto worker devices (Google Play distribution is not assumed).
\end{enumerate}


% -------------------------------------------------------
% APPENDICES
% -------------------------------------------------------
\newpage
\appendix

\section{Glossary}
\label{app:glossary}

\begin{description}[leftmargin=4cm, style=nextline]
    \item[API]
        Application Programming Interface. The defined set of endpoints
        through which the Android app and web dashboard communicate with
        the backend server.
    \item[APK]
        Android Package Kit. The installation file format for Android
        applications.
    \item[Audit Log]
        An immutable, chronological record of all privileged actions
        performed by administrators, including attendance edits, leave
        markings, petty-cash adjustments, and configuration changes.
    \item[Clock-In]
        The action of recording the start of a field worker's attendance
        session, captured with GPS coordinates and a selfie photograph.
    \item[Clock-Out]
        The action of recording the end of a field worker's attendance
        session.
    \item[CrashReporter]
        The interface through which the Android app reports unhandled
        exceptions and diagnostics. Firebase Crashlytics is the default
        implementation; it may be replaced or disabled via configuration.
    \item[Daily Rate]
        The configured gross wage amount earned by a worker for one full
        paid day of attendance. Used by the \texttt{DAILY\_FIXED} wage
        strategy. Currency is configurable per organisation.
    \item[Deployer]
        A technical operator who installs and maintains a self-hosted
        Beacon instance on behalf of an organisation.
    \item[FCM]
        Firebase Cloud Messaging. Google's cross-platform push
        notification service. The default \texttt{NotificationProvider}
        implementation.
    \item[Gazetted Holiday]
        A public holiday officially notified in a government gazette.
        The reference locale preset uses the Government of India gazette.
    \item[Geofence]
        A virtual geographic boundary defined by GPS coordinates (radius
        or polygon) used to determine whether a worker is within an
        acceptable work location.
    \item[GPS Accuracy]
        The estimated horizontal error radius (in metres) reported by
        the device GPS receiver at the time of a location reading.
        A lower value indicates a more precise fix.
    \item[JWT]
        JSON Web Token. A compact, URL-safe means of representing claims
        between two parties, used here for API authentication.
    \item[Locale Preset]
        A named configuration bundle containing a default holiday list,
        timezone, and currency for a specific deployment region.
        The India preset is shipped with the system; others may be
        contributed by the community.
    \item[NotificationProvider]
        The interface through which the backend dispatches push
        notifications to worker devices. Firebase Cloud Messaging is
        the default implementation.
    \item[Offline-First]
        An architectural pattern where all core application functions
        operate without internet connectivity, with data synced to the
        server when a connection becomes available.
    \item[Paid Leave]
        An absent working day explicitly marked by an administrator as
        paid; the day counts toward the worker's wage calculation.
    \item[Petty Cash]
        A small fund disbursed to workers for incidental work-related
        expenses. Managed via the petty-cash ledger feature.
    \item[Required Work Duration]
        The minimum number of hours per day configured by the
        administrator that constitutes a full-attendance session.
        Sessions shorter than this trigger a short-hours warning.
    \item[Short-Hours Flag]
        An indicator applied to an attendance record where the worker's
        session duration was less than the configured required work
        duration. Informational only; no automated penalty is applied.
    \item[Sync]
        The process of transmitting locally stored records (attendance,
        petty-cash expenses, configuration updates) between the Android
        device and the backend server.
    \item[TLS]
        Transport Layer Security. A cryptographic protocol providing
        secure communication over a network.
    \item[Unpaid Leave]
        An absent working day classified (explicitly or by default) as
        unpaid; the day is excluded from the worker's wage calculation.
    \item[UTC]
        Coordinated Universal Time. The primary time standard used for
        all timestamps in the system.
    \item[Wage Strategy]
        A named algorithm used to calculate a worker's gross wages.
        Built-in strategies: \texttt{DAILY\_FIXED},
        \texttt{MONTHLY\_DIVIDED}, \texttt{HOURLY}. Custom strategies
        may be registered by deployers.
    \item[Weekly Off]
        The day(s) of the week on which a worker is not expected to
        attend. Configured per worker; defaults to Sunday.
    \item[WorkManager]
        An Android Jetpack library for scheduling deferrable background
        tasks, used here for reliable sync and notification scheduling.
\end{description}
\newpage


\section{Issues List}
\label{app:issues}

The following open questions are pending resolution. Items marked with
$\star$ are particularly relevant to contributors picking up work on the
affected feature.

\begin{enumerate}
    \item \textbf{Clock-Out Selfie:} Should the clock-out flow also
          require a selfie photograph, or is GPS alone sufficient?
    \item \textbf{Edit Permissions:} Should both managers and admins be
          able to edit attendance records, or should this be restricted
          to admins only?
    \item $\star$ \textbf{Poor GPS Fallback:} What is the intended
          behaviour when GPS accuracy is poor (e.g., 100~m error radius)
          but the worker is visibly on-site? Options: (a) reject and
          require retry, (b) allow with a flag, (c) allow manager
          override remotely.
    \item \textbf{Half-Day Marking:} Is half-day attendance marking
          required in Version 1.0, or deferred to Phase 2? (See
          contribution roadmap, \S1.4.)
    \item \textbf{Camera Direction:} Should selfie capture use the front
          camera (true selfie) or the rear camera (higher quality)?
    \item \textbf{Offline Login Window:} What is the maximum acceptable
          duration for cached offline login before re-authentication
          is enforced?
    \item \textbf{Wage Precision:} Should rates support sub-unit decimal
          precision (e.g., paise for INR), or are whole-unit rates
          sufficient?
    \item \textbf{Wage Visibility on App:} Should the worker's wage
          summary be visible at all times on the home screen, or only
          on a dedicated earnings screen (privacy consideration)?
    \item \textbf{Holiday List Maintenance:} Who is responsible for
          updating the locale preset holiday list annually: the deploying
          organisation, or will the system provide an automated update
          mechanism?
    \item \textbf{Multiple Weekly Off Days:} Can a worker have more than
          one weekly off day (e.g., Saturday and Sunday)? Is a half-day
          weekly off pattern required?
    \item $\star$ \textbf{Petty-Cash Negative Balance:} Should the
          system allow a worker's petty-cash balance to go negative
          (i.e., can workers log expenses exceeding their current
          balance)?
    \item $\star$ \textbf{Petty-Cash Offline Expense:} If a worker adds
          an expense offline and the balance would go negative upon sync,
          how should the conflict be resolved?
    \item \textbf{Reminder Notification Time:} Is per-organisation
          configurability of reminder time sufficient, or should it also
          be configurable per worker?
    \item \textbf{Short-Hours Threshold Uniformity:} Is the required
          work duration the same for all workers, or should it be
          configurable per worker or per role?
    \item $\star$ \textbf{Multiple Geofence Zones:} Should a worker be
          permitted to clock in from any active geofence zone, or should
          zone assignment be per-worker?
    \item $\star$ \textbf{Alternative Notification Providers:}
          Which providers (ntfy, Gotify, APNS, SMTP) should be
          implemented for the v1.0 release, and which deferred to
          community contributions?
    \item \textbf{Custom Wage Strategies:} What interface (plugin
          contract) should a deployer implement to register a custom
          wage strategy? This needs a design decision before FR-13 is
          implemented.
\end{enumerate}

\end{document}
